\chapter{Tab 6 · Quality \& forensic review}\label{ch:review}

\section{Purpose}
The review tab answers file-integrity, result-consistency and control-identity questions. Its export contains plate map,
sample comparison, statistics, control charts, runs and QC, amplification curves and instrument records.
\readingnote{Use the sheet that answers the current check; the export does not replace the source instrument file.}

\figui{ui_review.png}{Review overview for the synthetic data set: experiments, stored and linked results, structural errors (the two planted errors), review findings.}{ui_review}

\section{The forensic log}
Every check writes events with a severity (\emph{error}, \emph{review}, \emph{info}), an area, a finding and the evidence.
The log collects structural checks, SOP findings and interpretation evidence. Control-chart alarm points are shown in Tab~2
and remain in the chart CSV/SVG exports; they are not duplicated here.

\par\medskip\noindent{\small\textbf{Listing}~\texttt{reviewForensicEvents()}\hfill\textcolor{gray}{HTML lines 4064--4071}}\par\nopagebreak
\lstinputlisting[style=vstyle,language=JavaScript,firstnumber=4064]{code/reviewForensicEvents.js}


\par\medskip\noindent{\small\textbf{Listing}~\texttt{reviewForensicEventsCoreUncached()}\hfill\textcolor{gray}{HTML lines 9068--9126}}\par\nopagebreak
\lstinputlisting[style=vstyle,language=JavaScript,firstnumber=9068]{code/reviewForensicEventsCore.js}


\par\medskip\noindent{\small\textbf{Listing}~\texttt{integrityEvents()}\hfill\textcolor{gray}{HTML lines 3762--3777}}\par\nopagebreak
\lstinputlisting[style=vstyle,language=JavaScript,firstnumber=3762]{code/integrityEvents.js}



\subsection*{Additional forensic checks}
\texttt{extraForensicEvents()} looks for identical curves or Cq values where independent data are expected (\emph{Copied data}),
dates in the wrong order, changes long after the run, stored calls that disagree with the SOP and similar signs.

\par\medskip\noindent{\small\textbf{Listing}~\texttt{extraForensicEventsUncached()}\hfill\textcolor{gray}{HTML lines 4659--4701}}\par\nopagebreak
\lstinputlisting[style=vstyle,language=JavaScript,firstnumber=4659]{code/extraForensicEvents.js}


\par\medskip\noindent{\small\textbf{Listing}~\texttt{westgardReview()}\hfill\textcolor{gray}{HTML lines 9232--9249}}\par\nopagebreak
\lstinputlisting[style=vstyle,language=JavaScript,firstnumber=9232]{code/westgardReview.js}



Exact control names are now the default. The SC profile adds a canonical-family view for descriptive pooling while retaining
the original sample text. Review baselines use the initial configured run count; pooling does not change stored values.

\section{Findings on the synthetic data}
{\footnotesize
\begin{longtable}{@{}lll@{}}
\toprule \textbf{Column 1} & \textbf{Column 2} & \textbf{Column 3} \\\midrule\endfirsthead
\toprule \textbf{Column 1} & \textbf{Column 2} & \textbf{Column 3} \\\midrule\endhead
error & Copied data & 1 \\\hline
review & Analysis membership & 3684 \\\hline
review & Control chart B-H3c & 122 \\\hline
review & SOP: DEMO laboratory SOP (synthetic example) & 102 \\\hline
review & Control chart B-H1 & 40 \\\hline
review & Detection format & 38 \\\hline
review & Control chart A-H4 & 35 \\\hline
review & Control chart B-H15 & 32 \\\hline
review & Control chart B-H11 & 31 \\\hline
review & Control chart B-S9 & 31 \\\hline
review & Control chart B-O1 & 31 \\\hline
review & Control chart B-H3 & 30 \\\hline
review & Control chart B-H3b & 30 \\\hline
review & Control chart B-H4 & 30 \\\hline
review & Control chart A-M1 & 27 \\\hline
review & Control chart B-H2 & 22 \\\hline
review & Timeline & 20 \\\hline
review & Provenance & 19 \\\hline
review & Control chart A-S4 & 13 \\\hline
review & Control chart B-H10b & 11 \\\hline
review & Control chart B-H5 & 10 \\\hline
review & Control chart B-H10 & 10 \\\hline
\bottomrule
\end{longtable}}


{\footnotesize
\begin{longtable}{@{}llll@{}}
\toprule \textbf{Experiment} & \textbf{Severity} & \textbf{Area} & \textbf{Finding and evidence} \\\midrule\endfirsthead
\toprule \textbf{Experiment} & \textbf{Severity} & \textbf{Area} & \textbf{Finding and evidence} \\\midrule\endhead
\texttt{DEMO\_LC1\_SCREEN\_013} & error & Copied data & 1 well pair(s) carry identical… \\\hline
\texttt{DEMO\_LC1\_SCREEN\_018} & review & Timeline & Last recorded change 410.0 h a… \\\hline
\texttt{DEMO\_LC1\_SCREEN\_018} & review & SOP: DEMO laboratory SOP (synt… & Control result — Reference che… \\\hline
\texttt{DEMO\_LC1\_SCREEN\_018} & review & SOP: DEMO laboratory SOP (synt… & Edited within the allowed wind… \\\hline
\texttt{DEMO\_LC1\_SCREEN\_018} & review & SOP: DEMO laboratory SOP (synt… & 15 result(s) need action (Repe… \\\hline
\bottomrule
\end{longtable}}



The synthetic gallery includes a dilution-series example in which one NTC shows a late signal; the example demonstrates a
failed negative-control rule without reproducing a laboratory result.

Most review findings are expected for this data set: many wells are named but not included in an analysis
(\emph{Analysis membership}), and the history contains the planted trends, which the control charts report as alarms.
\readingnote{A red point is outside a displayed limit; an amber point is an enabled pattern rule inside the limits. Both
are review prompts. A missing chart series means that its data or eligibility requirements were not available.}

\fig[0.9]{g_all_x_findings.png}{\texttt{x\_findings} — Pareto of error and review findings by area.}{findings}


